\documentclass[11pt]{article}
\usepackage[margin=1.1in]{geometry}
\usepackage{amsmath,amssymb,mathtools,booktabs}
\usepackage[colorlinks=true,linkcolor=blue,citecolor=blue,urlcolor=blue]{hyperref}

\newcommand{\Mone}{M_1}
\newcommand{\Mtwo}{M_2}
\newcommand{\what}[1]{\widehat{#1}}

\title{Noise-whitened VarPro fitting: combining smooth and discrete\\
       basis functions without mass matrices in the fit}
\author{N.~Alger}
\date{}

\begin{document}
\maketitle

\begin{abstract}
\noindent
Per-row VarPro fitting of the LG-PSF Hessian model combines two kinds of basis
function: theta-dependent smooth (Laguerre--Gaussian) features, and a fixed
theta-independent ``extra'' basis (a discrete spike, and generalizations of
it) supplied directly as sampled values. The two live in different spaces
with different natural inner products (one in the discretized function space,
one in its dual), and combining them naively -- e.g.\ orthogonalizing them
against each other under a plain Euclidean inner product -- either gets the
wrong answer or requires the fitting code to know about the mesh mass
matrices explicitly. This note derives a \emph{noise-whitening} change of
variables under which the entire per-row fit becomes an ordinary
(mass-free) least-squares problem: every basis function, every derivative,
and the data itself are rescaled once, at the interface, so that plain
Euclidean operations throughout the fit are automatically correct. The two
mass matrices ($M_1$ for the row/target space, $M_2$ for the column/source
space) appear only in that one whitening layer; the VarPro fitting machinery
itself never sees them. The whitening transform is theta-independent and
symmetric, so it composes with the existing forward/reverse-mode derivative
machinery (\texttt{lg\_functions.hpp}, \texttt{ellipsoid\_transform.hpp}) with
no new derivative math required.
\end{abstract}

\section{Setup}

Recall the discretization picture: $X$ is the discrete
source/column function space with inner product $\langle u,v\rangle_X =
u^\top \Mtwo v$; $Y$ is the discrete target/row space with $\langle
u,v\rangle_Y = u^\top \Mone v$; $X'$, $Y'$ are their duals, with $\Mtwo:
X\to X'$ and $\Mone: Y\to Y'$ the Riesz maps induced by those inner
products; the kernel $\Phi: X'\to Y$ and the assembled operator $H = \Mone
\Phi \Mtwo : X \to Y'$. Both $\Mone,\Mtwo$ are diagonal (lumped mass).

Each Laguerre--Gaussian smooth feature $\phi_i(\cdot;\theta)$, sampled at
mesh nodes, is an element of $X$ -- and, because the continuum family
$\{\psi_i\}$ is exactly $L^2(\mathbb{R}^N)$-orthonormal (verified numerically
in \texttt{test\_lg\_functions.cpp} by tensor-product Gauss--Hermite
quadrature, which pins the generated harmonic table, the Laguerre recurrence
and the normalization at once), the
sampled $\phi_i$ are nearly $\Mtwo$-orthonormal: $\phi_i^\top \Mtwo \phi_{i'}
\approx \delta_{ii'}$, up to mesh resolution. A user-supplied ``extra''
basis function $e_d$ (a one-hot vector for the diagonal spike; more general
indicator combinations for e.g.\ a ring-neighbor correction) is instead a
sampled \emph{functional} -- the coordinate representation of a discrete
measure (a sum of point masses) via the FEM nodal interpolation property
$N_k(x_i) = \delta_{ik}$ -- and so is native to $X'$, not $X$.

\section{The row model}

Fix a row $\rho$ with mass $m_\rho := (\Mone)_{\rho\rho}$, and a local
neighbor/support batch of $K$ column points $x_j$ with per-point masses
$m_j := (\Mtwo)_{jj}$. The row model is
\begin{equation}
H[\rho,j] \;\approx\; \underbrace{m_\rho m_j \sum_i c_i\, \phi_i(x_j;\theta)}_{\text{smooth}}
\;+\; \underbrace{m_\rho \sum_d s_d\, e_d(j)}_{\text{extra}},
\label{eq:row-model}
\end{equation}
with $c_i$, $s_d$ the (linear) coefficients VarPro solves for at each trial
$\theta$. The column mass $m_j$ on the smooth term is a quadrature weight --
$\phi_i$ approximates a continuum kernel evaluated at a mesh quadrature
point, and $m_j$ discretizes the integral. The extra term carries no
quadrature weight because it is not a continuum-kernel evaluation at all: it
is a direct, discrete correction to specific matrix entries. This asymmetry
is what reproduces the $m_\rho s_\rho\,\delta_{\rho j}$ spike convention
exactly.

Probing: draw i.i.d.\ standard normal $z_\ell$ ($\ell = 1,\dots,k$,
restricted to the batch -- a restriction of an i.i.d.\ vector is still
i.i.d.), observe $y_\ell(\rho) = (Hz_\ell)(\rho) = \sum_j H[\rho,j]\,
z_\ell(j)$. Substituting \eqref{eq:row-model},
\begin{equation}
y_\ell(\rho) \;=\; m_\rho \sum_i c_i\, g_\ell^{(i)} \;+\; m_\rho \sum_d s_d\, h_\ell^{(d)},
\qquad
g_\ell^{(i)} := (\Mtwo \phi_i)\cdot z_\ell, \quad h_\ell^{(d)} := e_d \cdot z_\ell.
\label{eq:probe-eq}
\end{equation}

\paragraph{The problem with combining these naively.} $\phi_i$ (in $X$) and
$e_d$ (in $X'$) are different kinds of object; comparing or orthogonalizing
them requires first mapping one into the other's space via the Riesz map,
and then using \emph{that} space's natural inner product. Using the raw
Euclidean inner product throughout happens to reproduce the correct answer
for a single one-hot $e_d$ -- every
diagonal-respecting inner product agrees on a coordinate-axis-aligned vector
-- but is not correct in general (e.g.\ a ring-neighbor indicator combining
several points), and more importantly requires the fitting code to reach
into $\Mtwo$ (and its inverse) explicitly to do the orthogonalization
correctly. That is exactly the kind of mass-matrix leakage this note is
trying to eliminate.

\section{Noise whitening}

Define, per row (all quantities below implicitly depend on $\rho$ through
$m_\rho$ and the batch):
\begin{align}
\what{z}_\ell &:= \Mtwo^{1/2} z_\ell, \label{eq:whiten-z}\\
\what{\phi}_i(\theta) &:= \sqrt{m_\rho}\; \Mtwo^{1/2} \phi_i(\theta), \label{eq:whiten-phi}\\
\what{e}_d &:= \sqrt{m_\rho}\; \Mtwo^{-1/2} e_d, \label{eq:whiten-e}\\
\what{y}_\ell &:= y_\ell(\rho) / \sqrt{m_\rho}. \label{eq:whiten-y}
\end{align}
$\Mtwo^{1/2}$ is diagonal, hence symmetric, so $(\Mtwo \phi_i)\cdot z_\ell =
(\Mtwo^{1/2}\phi_i)\cdot(\Mtwo^{1/2}z_\ell)$, giving $\what\phi_i \cdot
\what z_\ell = \sqrt{m_\rho}\, g_\ell^{(i)}$; likewise $\what e_d \cdot \what
z_\ell = \sqrt{m_\rho}\,(\Mtwo^{-1/2}e_d)\cdot(\Mtwo^{1/2}z_\ell) =
\sqrt{m_\rho}\, e_d\cdot z_\ell = \sqrt{m_\rho}\, h_\ell^{(d)}$. Dividing
\eqref{eq:probe-eq} by $\sqrt{m_\rho}$:
\begin{equation}
\what y_\ell \;\approx\; \sum_i c_i \big(\what\phi_i(\theta)\cdot\what z_\ell\big)
\;+\; \sum_d s_d \big(\what e_d \cdot \what z_\ell\big).
\label{eq:whitened-model}
\end{equation}
This is an ordinary (unweighted) linear regression in $\ell$: every mass
matrix has been absorbed into the four definitions
\eqref{eq:whiten-z}--\eqref{eq:whiten-y}, and \eqref{eq:whitened-model}
itself is mass-free.

\emph{Note on $\what e_d$:} $\Mone$ contributes the same $\sqrt{m_\rho}$
factor to \emph{both} terms of \eqref{eq:row-model} (there is one row mass,
shared), so $\what e_d$ needs the same $\sqrt{m_\rho}$ prefactor as $\what
\phi_i$ even though it does not involve $\Mtwo$ (which enters with the
opposite sign of exponent, $-1/2$ vs.\ $+1/2$, reflecting that $e_d$ lives in
$X'$ where $\phi_i$ lives in $X$). This factor is easy to drop, and dropping
it is invisible to every check that does not touch $\Mone$ or the extra basis
-- finite differences and adjoint consistency on the feature layer both pass
without it. \texttt{test\_whitening.cpp} guards it directly, by constructing an
explicit $H$ from \eqref{eq:row-model} and checking that the whitened
quantities reproduce $y_\ell(\rho)$ to machine precision.

\section{Consequences}

\paragraph{Orthonormality is inherited directly.} $\what\phi_i\cdot\what
\phi_{i'} = m_\rho\,(\Mtwo^{1/2}\phi_i)\cdot(\Mtwo^{1/2}\phi_{i'}) = m_\rho\,
\phi_i^\top \Mtwo \phi_{i'} \approx m_\rho\,\delta_{ii'}$: the whitened
smooth basis vectors are, in the ordinary Euclidean sense, exactly as
orthonormal (up to the constant $m_\rho$) as the $\Mtwo$-orthonormality
already established for the sampled LG modes. No separate weighted inner
product is needed to see this -- it falls out of using the plain dot product
on already-whitened vectors.

\paragraph{Plain-Euclidean orthogonalization is now correct by construction.}
Projecting a smooth feature against an extra basis correctly requires the
$\Mtwo^{-1}$-weighted inner
product on $X'$: $\what\phi_p^\flat \mapsto \what\phi_p^\flat -
e_i\,(e_i^\top \Mtwo^{-1} e_i)^{-1}(e_i^\top \Mtwo^{-1} \what\phi_p^\flat)$,
$\what\phi_p^\flat := \Mtwo\phi_p$. Checking this against the whitened
quantities here: $\what\phi_p \cdot \what e_i = (\Mtwo^{1/2}\phi_p)\cdot
(\Mtwo^{-1/2}e_i) = \phi_p \cdot e_i$ and $\what e_i \cdot \what e_i = e_i^\top
\Mtwo^{-1} e_i$ (the $\sqrt{m_\rho}$ factors cancel in the ratio) --
identical projection coefficients. So plain-Euclidean orthogonalization of
$\what\phi_p$ against $\what e_i$ \emph{is} the earlier, carefully-derived
$\Mtwo^{-1}$-weighted projection; whitening does not introduce a different
answer, it removes the need to ever justify or implement a non-Euclidean
inner product near the fitting code.

\paragraph{Derivatives compose for free.} $\sqrt{m_\rho}\,\Mtwo^{1/2}$ and
$\sqrt{m_\rho}\,\Mtwo^{-1/2}$ are fixed, $\theta$-independent, symmetric
linear operators (diagonal matrices). Consequently:
\begin{itemize}
\item the JVP of $\what\phi_i$ is that operator applied to the JVP of the
  raw $\phi_i(\theta)$ -- i.e.\ apply it to whatever
  \texttt{grad\_eval\_lg\_nd} $\cdot$ \texttt{jvp\_T} (chain rule through
  the pullback) already produces;
\item because the operator is symmetric, the VJP is the \emph{same}
  operator applied to the incoming cotangent \emph{before} it is fed into
  the existing VJP chain.
\end{itemize}
This is the same ``compose with a fixed linear map'' pattern that
\texttt{ellipsoid\_transform.hpp} already uses to chain its two stages. No new
derivative formulas are needed anywhere.

\section{Resulting architecture}

\begin{itemize}
\item \textbf{Mass-free composed feature layer} (\texttt{lg\_ellipsoid\_feature.hpp}):
  $\phi_i(x;\theta) := \psi_i(T(\theta,x))$ and its JVP/VJP, built purely by
  composing the two existing modules via the chain rule. Knows nothing
  about mass matrices.
\item \textbf{Whitening interface layer} (\texttt{whitening.hpp}): the only
  place $\Mone,\Mtwo$ appear. Provides (a) \texttt{whiten\_probes},
  \texttt{whiten\_data}, \texttt{whiten\_extra} for the user-supplied raw
  probe/data/extra-basis arrays, and (b) whitened wrappers around the
  feature layer's eval/JVP/VJP, applying
  $\sqrt{m_\rho}\,\Mtwo^{\pm 1/2}$ once.
\item \textbf{Generic VarPro core} (\texttt{varpro.hpp}): takes whitened
  probes, whitened data, and a basis functor $\theta \mapsto$ an evaluation
  offering values and a VJP. Performs the per-row fit in ordinary Euclidean
  terms throughout and never references a mass matrix. Concretely it projects
  the extra block out once by Frisch--Waugh--Lovell rather than at every trial
  point, solves the inner least-squares problem, forms the reduced residual,
  and hands its $\theta$-Jacobian to a trust-region Levenberg--Marquardt loop
  (\texttt{detail/levenberg\_marquardt.hpp}).
\end{itemize}

The prediction that this layering would hold has since been tested by building
on top of it. Two things are worth recording because they confirm it.
% (operator-level QC section added below, 2026-08-13 -- draft for review)

First, the fitting core is genuinely generic: \texttt{varpro.hpp} is templated
on the basis functor and contains nothing specific to Laguerre--Gaussians,
ellipsoids or point-spread functions. It is usable for any separable
least-squares problem.

Second, the whitening layer stayed exactly one file. Every mass matrix in the
library still appears only in \texttt{whitening.hpp}, and everything above it
--- the candidate search, the mode ladder, the cross-validation scoring, the
whole-operator layer --- works in whitened coordinates without ever needing
$\Mone$ or $\Mtwo$. The one place a mass reaches higher is a scalar: the
operator layer passes each row's own $m_\rho$ down as \texttt{target\_mass}, so
the returned coefficients come back in physical units. Because the whitened
design matrix is column-equilibrated inside the inner solve, that scalar
rescales only $(c, s)$ --- the fitted $\theta$, the scores and the entire
selection are invariant to it.

\section{Operator-level quality control}
\label{sec:operator-qc}

The fit makes two different kinds of decision, and they deserve two different
metrics. The per-row cross-validation score selects a \emph{model for one row}
(how many LG modes to use), and it already operates in whitened coordinates:
its residuals are $\what y$-residuals, i.e.\ residuals in the $y/\sqrt{m_\rho}$
scaling of \eqref{eq:whiten-y}. Nothing in this section touches it. The second
decision is made about the \emph{operator as a whole} -- is the assembled fit
$B \approx H$ accurate enough to stop spending probes? -- and this section
derives the metric for that decision, extending the whitening story from rows
up to the operator.

\paragraph{Why a naive global metric misleads.} A natural-looking choice is to
hold out a few probes and average their relative residuals,
$\tfrac1k\sum_\ell \|Bz_\ell - y_\ell\| / \|y_\ell\|$. The trouble is the
denominator: $\|Hz\|$ is dominated by the overlap of $z$ with the leading
spectral directions of $H$, and the operators we care about have large spectral
outliers -- that is the very reason they need preconditioning. Each per-probe
ratio therefore has a heavy-tailed denominator, and with a handful of held-out
probes the \emph{mean of ratios} inherits the tail: its value is largely a
property of the draw, not of the fit. (In the ice-sheet application we observed
a factor-of-three spread across probe draws at fixed fit quality.) The fix is
the same one used throughout these notes: state the quantity in the natural
norms, whiten, and then let plain Euclidean arithmetic be correct by
construction.

\paragraph{The metric.} The error $D := H - B$ maps $X \to Y'$, so its natural
Hilbert--Schmidt norm measures inputs in $\langle u,v\rangle_X = u^\top \Mtwo
v$ and outputs in the dual norm $\|r\|_{Y'}^2 = r^\top \Mone^{-1} r$. Summing
$\|D u_k\|_{Y'}^2$ over an $X$-orthonormal basis $\{u_k\}$ gives
\begin{equation}
\|D\|_{\mathrm{HS}(X,Y')}^2
\;=\; \big\| \Mone^{-1/2}\, (H - B)\, \Mtwo^{-1/2} \big\|_F^2 .
\label{eq:hs-norm}
\end{equation}
That is, the natural operator error is the plain Frobenius norm of the error
after whitening both sides -- the operator-level analogue of the per-row
transforms \eqref{eq:whiten-z}--\eqref{eq:whiten-y}, with the single row mass
$m_\rho$ replaced by the full $\Mone$. The quality-control quantity is the
relative version, $\|D\|_{\mathrm{HS}} / \|H\|_{\mathrm{HS}}$.

\paragraph{The estimator.} Draw held-out probes i.i.d.\ standard normal
\emph{in the whitened variables}: $\what z_\ell \sim N(0, I)$, i.e.\ $z_\ell =
\Mtwo^{-1/2} \what z_\ell$ in raw coordinates, so that $\mathbb{E}\, z_\ell
z_\ell^\top = \Mtwo^{-1}$. In function-space terms this makes $z_\ell$ white
noise with respect to the $X$ inner product -- coefficient covariance
$\Mtwo^{-1}$ is the FEM representation of $L^2$-white noise -- whereas a
coordinate-i.i.d.\ draw sits one factor of $\Mtwo$ away from it. Then
\begin{equation}
\mathbb{E}\, \big\| \Mone^{-1/2} (H-B) z_\ell \big\|^2
\;=\; \|D\|_{\mathrm{HS}(X,Y')}^2 ,
\qquad
\mathbb{E}\, \big\| \Mone^{-1/2} H z_\ell \big\|^2
\;=\; \|H\|_{\mathrm{HS}(X,Y')}^2 ,
\end{equation}
and the quality control is the ratio of sums (an \emph{energy} ratio, not a
mean of ratios):
\begin{equation}
\mathrm{qc}^2 \;=\;
\frac{\sum_\ell \big\| \Mone^{-1/2} (B z_\ell - y_\ell) \big\|^2}
     {\sum_\ell \big\| \Mone^{-1/2} y_\ell \big\|^2},
\qquad y_\ell = H z_\ell .
\label{eq:qc-energy}
\end{equation}
Two remarks on this estimator. First, as a ratio of sums it is biased at
$O(1/k)$ for the ratio of expectations, which is immaterial for a stopping
test. Second, and decisively, it concentrates where the mean of ratios does
not: numerator and denominator fluctuate \emph{together} through the shared
draw, so their ratio is far more stable than either sum alone. The heavy tail
that poisons the per-probe ratios enters both sums the same way and largely
cancels.

\paragraph{Held-out interpretation without the i.i.d.\ assumption.} The
i.i.d.-in-whitened-variables assumption is what upgrades \eqref{eq:qc-energy}
to an unbiased estimate of the \emph{global} relative Hilbert--Schmidt error;
it is not what makes the formula meaningful. For any held-out probe family --
scaled or unscaled, random or structured (polynomials, sinusoids, Dirac
combs, draws from other smoothness classes) -- \eqref{eq:qc-energy} is still
a perfectly reasonable held-out relative error: the error of $B$'s action on
that family, measured in the natural norms. What changes with the probe
family is only \emph{which} operator error the number reflects -- the family
weights the directions being tested -- not whether the number is a valid
relative error on them.

\paragraph{Mesh independence.} Whitening is what makes the number mean the same
thing on different meshes. Heuristically, \eqref{eq:hs-norm} is the quadrature
approximation of the continuum Hilbert--Schmidt (relative) norm of the kernel
error, so under refinement -- uniform or adaptive -- the metric approaches a
mesh-independent limit whenever the underlying kernel is square-integrable. We
do not prove this; it is the standard quadrature heuristic. By contrast, the
unweighted version (Euclidean norms, coordinate-i.i.d.\ probes) estimates the
raw-coordinate Frobenius ratio, which weights every node equally: on an
adaptive mesh, refining a region silently re-weights the metric toward it. On
a quasi-uniform mesh the two versions differ only by constants that cancel in
the ratio, which is why the distinction is easy to miss until an adaptive mesh
makes it matter.

\paragraph{What changes and what does not.} In whitened variables the probing
rule is unchanged -- ``draw $\what z$ i.i.d.'' is exactly the
\eqref{eq:whiten-z} convention -- and the residual weighting $\Mone^{-1/2}$ is
the global form of the per-row $y/\sqrt{m_\rho}$. The per-row scores, the mode
selection, and the fitted operator are untouched. Two consequences do warrant
attention. First, in raw coordinates the probes acquire a factor
$\Mtwo^{-1/2}$, and the fit consumes the same probes as the quality control:
changing the probe covariance changes the fitted operator (incidentally, it
makes the whitened per-row design itself i.i.d., since
$\what z = \Mtwo^{1/2} z = $ the drawn vector), so adopting
\eqref{eq:qc-energy} is a change to the fit's inputs, not only to a printed
number. The per-row fits are insensitive to this change: a row's regression
is valid conditioned on whatever probes were drawn -- the probe distribution
affects statistical efficiency, not correctness -- and within one PSF window
the column masses vary mildly, so the covariance change amounts to a
near-constant factor per window, which the regression absorbs. It is only
the \emph{global} quantity, on a mesh whose resolution varies by orders of
magnitude across the domain, that needs the scaling. Second, any threshold
calibrated against a previous metric must be recalibrated against this one,
once.

\end{document}
